\section{Ergebnisse}

\subsection{Repräsentative End-to-End-Artefaktkette}

Der vollständige Funktionsnachweis wird anhand eines intern konsistenten
Referenzlaufs dargestellt. Die endgültige Fassung zeigt dafür dieselbe
Verarbeitungskette in aufeinander abgestimmten Ansichten: Rohframe,
\texttt{default}/\texttt{middle}-Maske, ideale Bild-/Maskendomäne,
COLMAP-Sparse-Modell, objektbezogene STS-Gaussians, Route-A-Coarse-Mesh sowie
lokale Centerline und Grad-10-B-Spline. Eine solche Kette ist für alle drei
Auflösungsstufen visuell dokumentiert; Anhang~\ref{app:durchlauf} zeigt den
vollständigen Pipeline-Durchlauf mit 5 FPS und OPENCV in 720p, qHD und Low.

Die archivierten Läufe verwenden durchgängig 240 Kameras; davon dienen 210 dem
Training und 30 der Evaluation. Der feste Eval-Split und seine Begründung sind
in Kapitel~\ref{sec:versuchsaufbau} beschrieben; hier genügt die Erinnerung,
dass die Evaluationsbilder gerade deshalb bewertbar sind, weil sie nicht am
Training beteiligt waren. Jede archivierte Kette erzeugte Route-A-Mesh,
Rendering, objektmaskierte Metriken, Centerline, B-Spline und GeoJSON. Der
Autopilotnachweis selbst ist über sechs archivierte Autopilot-Volläufe erbracht
(Abschnitt~\ref{sec:versuchsaufbau}); die Visualisierung dieser Durchläufe in
ausgewählten Screenshots liegt ebenfalls im Anhang~\ref{app:durchlauf}. Die
historischen Zwischenstände, die zu den heutigen Einstellungen in STS,
Objektfilterung und SuGaR geführt haben, sind in
Kapitel~\ref{sec:implementierung} sowie in
Abschnitt~\ref{sec:sugar-folgepruefung} dokumentiert.

\subsection{Visueller Vergleich der Auflösungsvarianten}
\label{sec:visueller-aufloesungsvergleich}

Die visuelle Gegenüberstellung der drei Auflösungen ist in
Anhang~\ref{app:durchlauf} entlang derselben Verarbeitungskette dokumentiert.
Bei
den Masken ist in den höher aufgelösten Varianten eine feinere Kontur- und
Randlokalisierung erkennbar. In keiner der drei Varianten treten jedoch
offensichtliche gravierende Fehler wie eine vollständig leere Maske, ein
Verlust des Zielobjekts oder große Ausrisse auf. In der \texttt{low}-Variante
wird zusätzlich das Verbindungskabel zwischen den Rohren in die Objektmaske
aufgenommen. Streng semantisch ist dies eine Übersegmentierung; da das Kabel
direkt zum dargestellten Objektverbund gehört, wird dieser Anteil für die
qualitative Pipelinebewertung als vertretbar behandelt.

In den Übersichten der COLMAP-Punktwolken sind zwischen den Auflösungen keine
direkten, großräumigen Unterschiede zu erkennen. Dies ist mit dem festgelegten
SIFT-Limit von 4096 Merkmalen pro Bild vereinbar: Eine höhere Pixelauflösung
führt nicht automatisch zu einer proportional größeren Punktwolke. In der
vergrößerten Darstellung fällt jedoch auf, dass in der linken unteren
Wandecke der \texttt{720p}-Variante mehr Punkte sichtbar sind. Dies kann darauf
hindeuten, dass das begrenzte Merkmalsbudget dort anders und räumlich breiter
über die ansonsten gleichförmige Fläche verteilt wird. Eine breitere Verteilung
homologer Bildpunkte könnte die Track-Stabilität und damit die interne
Geometriekonsistenz des SfM-Modells begünstigen. Aus den Screenshots allein
folgt daraus jedoch kein Nachweis einer höheren absoluten geometrischen
Korrektheit; dafür wären insbesondere Trackstatistiken, räumliche
Korrespondenzverteilungen und eine unabhängige Referenz erforderlich.

Bei den Splats ist der auffälligste Unterschied in der \texttt{low}-Variante
die bereits genannte Aufnahme des Verbindungskabels. Außerdem unterscheiden
sich die oberen Objektkanten: In den höher aufgelösten Varianten wirkt die
Objektgeometrie über die Blickwinkel hinweg konsistenter aufgenommen, während
die \texttt{low}-Variante am unteren rechten Bildrand einen zusätzlichen
Floater zeigt. Die höher aufgelösten Splats wirken dadurch in dieser
Gegenüberstellung geometrisch sauberer. Auch diese Aussage ist als visueller
Befund und nicht als unabhängige Genauigkeitsmessung zu verstehen.

Die Meshbilder zeigen eine entsprechende Übertragung der Unterschiede, weil
die hier dargestellte Route~A das Mesh unmittelbar aus den STS-Splats ableitet.
In den Detailansichten weist die niedrigste Auflösung die ungenaueste und
zackigste Oberfläche mit mehreren Spitzen auf. Die \texttt{qHD}-Variante wirkt
in der betrachteten Ansicht überraschend glatt und insgesamt visuell am
geschlossensten. Die \texttt{720p}-Variante erscheint dagegen besonders klar
durch die Maskengrenze begrenzt: Aufgenommene Bereiche besitzen deutliche
Kanten, ohne dass ein starkes Glätten oder Ausfransen entsteht. Dieses
Ausfransen ist bei der \texttt{low}-Centerline am deutlichsten; sie enthält
mehr Richtungswechsel und Ecken. Die Centerlines von \texttt{720p} und
\texttt{qHD} unterscheiden sich visuell dagegen nur gering. Die Befunde sind
qualitativ und ersetzen keine Vermessung gegen eine unabhängige Centerline.

\subsection{COLMAP-Voruntersuchung}

Ebenso wurde vor den Matrixläufen eine Versuchsreihe mit COLMAP gestartet;
dabei wurde geprüft, inwiefern sich einzelne Einstellungen im Hinblick auf den
zeitlichen Aufwand und die Qualität der Sparse-Rekonstruktion lohnen. Alle
fünf dokumentierten Hauptläufe registrierten 240 von 240 Bildern; die
vollständige Tabelle steht in Anhang~\ref{tab:colmap-vorstudie}. Guided
Matching erhöhte die Sparse-Punktzahl gegenüber Plain-SIFT 4096 nur um rund
1,98\,\%, benötigte aber nahezu die doppelte Zeit und erhöhte den mittleren
Reprojektionsfehler. 8192 Merkmale lieferten etwa 0,82\,\% mehr Punkte bei
5,5\,\% mehr Laufzeit.

Diese Beobachtungen widersprechen einander nicht: Bei fixiertem
Merkmalslimit ist die Anzahl detektierter Merkmale pro Bild weitgehend
identisch, weshalb die Punktzahl kaum steigt -- mehr Eingabemerkmale führen
nicht automatisch zu mehr 3D-Punkten. Guided Matching ergänzt zusätzlich
Korrespondenzen an schwierigen Stellen; dass dabei der mittlere
Reprojektionsfehler leicht ansteigt, zeigt, dass diese zusätzlichen Matches
im Mittel schwerer zu konsistenten Spuren zu bündeln sind als das
Grundmuster. Beides zusammen rechtfertigt für den Projektstand den Standard
Plain-SIFT 4096 ohne Guided Matching: gleiche Registrierungsqualität bei
geringerer Laufzeit und einfacherer Fehlerdiagnose.

\subsection{Objektmaskierte Bildmetriken}

Ein Matrixlauf ist ein vollständig protokollierter Einzelversuch für genau
eine Kombination aus Kameramodell, FPS, Auflösung und Meshroute. Der Runner
erzeugt dafür einen isolierten Frame- und Maskensatz, führt COLMAP, die ideale
Bilddomäne, STS, die gewählte Meshroute, Rendering und objektmaskierte
Metriken aus und archiviert anschließend die Resultate. In der Entwicklung
traten unter anderem ein durch Unterprozesse konsumierter TSV-Standardinput,
leere feste Eval-Masken sowie ein SuGaR-Import-/Tensorlayoutfehler auf. Diese
Fehler wurden durch stdin-unabhängige Schleifen, einen nichtleeren festen
Eval-Split sowie die Korrektur des Renderhelfers behoben. Dadurch lassen sich
die Matrixläufe heute mit denselben Parametern erneut erzeugen und unter
identischen Kriterien vergleichen.

Die Matrix dient dabei nicht dazu, aus einer einzelnen Bildmetrik eine
„beste“ Pipeline zu küren. Sie zeigt, welche Kombinationen technisch
durchlaufen, welche Fehlerfälle auftreten und wie sich die Ansichtsmetriken
unter kontrollierten Änderungen verhalten. Die Produktionsentscheidung für
Route A stützt sich deshalb vor allem auf die A/B/C/D-Geometrieablation und
den direkten Erhalt der STS-Gaussian-Geometrie; die Matrix ergänzt diesen
Befund um Robustheit und reproduzierbare Vergleichswerte.

Die STS-Grafik verwendet ausschließlich \path{sts_masked.json}, die
SuGaR-Grafik \path{sugar_coarse_masked.json}; die Interpretation der drei
Ansichtsmetriken und ihre Grenzen wurden bereits in
Kapitel~\ref{sec:grundlagen} definiert und werden hier nicht wiederholt.
Für die Meshentscheidung sind zusätzlich Vertex-/Face-Struktur, relative
Meshabstände und die visuelle Artefaktprüfung maßgeblich.

Ein erneuter vollständiger Matrixlauf ist als Wiederholungsprüfung sinnvoll.
Bei nominal gleichen Parametern sind kleine Abweichungen durch stochastische
Optimierung, Surface-Sampling sowie Software- und Hardwarestände möglich. Eine
Abweichung wird erst dann als technisch relevant bewertet, wenn Commitstände,
Eval-Split, Framezahl, Maskenabdeckung und Seeds übereinstimmen. Für die PA
werden daher nicht alle wiederholten Arbeitskopien als eigenständige
wissenschaftliche Ergebnisse benötigt; entscheidend sind die standardisierten
Metriken, die Endprodukte und die reproduzierbaren Metadaten.

Die zentrale Metrikübersicht ist nun nach Modellstufe getrennt. Die STS-Grafiken
lesen ausschließlich vollständige Original-GS-Route-A-Läufe aus
\texttt{sts\_masked.json}; die SuGaR-Grafiken lesen ausschließlich
\texttt{sugar\_coarse\_masked.json}. Vorhandene STS-Dateien aus später
fehlgeschlagenen SuGaR-Routen bleiben als Fehlernachweis archiviert, werden aber
nicht zusätzlich in der STS-Auswertung dargestellt. Die beiden FPS-Profile und
die jeweilige Runauflösung bleiben sichtbar; ein globales Ranking über
verschiedene Auflösungen wird nicht vorgenommen.

\begin{figure}[H]
  \centering
  \includegraphics[width=\textwidth]{sts_masked_overview.pdf}
  \caption{Getrennte STS-7000-Auswertung der vollständigen Original-GS-Route A
  aus den archivierten \texttt{sts\_masked.json}-Dateien.}
  \label{fig:pa-matrix-overview}
\end{figure}

\begin{figure}[H]
  \centering
  \includegraphics[width=\textwidth]{sts_masked_per_frame_boxplots.pdf}
  \caption[STS-Einzelansichten]{Verteilung der Einzelansichten in den
  vollständigen Original-GS-Route-A-Läufen. Punkte entsprechen den in
  \texttt{per\_frame} archivierten Ansichtsmetriken. SR ist blau, OPENCV
  orange und PINHOLE grau dargestellt.}
  \label{fig:pa-sts-boxplots}
\end{figure}

\begin{figure}[H]
  \centering
  \includegraphics[width=\textwidth]{sugar_coarse_masked_overview.pdf}
  \caption{Getrennte SuGaR-Coarse-9000-Auswertung. Die STS-Baseline wird
  nicht als SuGaR-Wert verwendet.}
  \label{fig:pa-sugar-overview}
\end{figure}

\begin{figure}[H]
  \centering
  \includegraphics[width=\textwidth]{sugar_coarse_masked_per_frame_boxplots.pdf}
  \caption[SuGaR-Coarse-Einzelansichten]{Verteilung der Einzelansichten in
  den zwölf SuGaR-Coarse-Läufen. Punkte entsprechen den in
  \texttt{per\_frame} archivierten Ansichtsmetriken.}
  \label{fig:pa-sugar-boxplots}
\end{figure}

Die Boxplotpunkte liegen jeweils in der x-Kategorie ihres Laufs und werden nur
leicht horizontal verteilt; die y-Achsen berücksichtigen alle archivierten
Einzelwerte einschließlich Ausreißern.

Die deutlich niedrigeren SuGaR-Werte gegenüber den STS-Werten sind kein
direkter Widerspruch. Die beiden Dateien beschreiben unterschiedliche
Modellstufen: Der STS-Splat wird vor der zusätzlichen SuGaR-Optimierung und dem
SuGaR-Rendering ausgewertet. Im vollständigen Repeat-Lauf beträgt der
beispielhafte 5-FPS-720p-Wert für SIMPLE\_RADIAL 29,43 dB PSNR, 0,897 SSIM und
0,122 LPIPS im STS-Splat; der zugehörige SuGaR-Coarse-Wert beträgt 21,22 dB,
0,778 SSIM und 0,151 LPIPS. Die Repeat-Grafiken verwenden diese getrennten
Metrikdateien fachlich korrekt. Ältere STS-Baselines aus fehlgeschlagenen
SuGaR-Routen bleiben historisch nachvollziehbar, werden aber nicht als
SuGaR-Ergebnisse bezeichnet.

In der jeweiligen Runauflösung zeigt OPENCV-A-Low bei 5 FPS besonders hohe
Bildmetriken. Bei SIMPLE\_RADIAL und OPENCV bezieht sich die Auswertung dabei
auf die ideale, gewarpte Bilddomäne; „jeweilige Runauflösung“ bezeichnet nur
die Pixelauflösung innerhalb dieser Domäne. Dieser Befund ist ein
Screening-Ergebnis. Niedrige Auflösung kann Kantenversatz und hochfrequente
Fehler glätten, sodass bessere Bildmetriken nicht automatisch eine bessere
3D-Geometrie bedeuten. Die Werte werden daher nicht als alleinige
Produktionsentscheidung verwendet.

Die Metrikwerte verhalten sich nicht monoton über die Konfigurationen: 5 FPS
liefern nicht in jeder Kombination den numerisch besten Wert, und Low-Auflösung
ist nicht in jedem Modell besser als qHD. Die beobachteten Unterschiede
entstehen daher wahrscheinlich aus einer Interaktion von Bildauflösung,
zeitlicher Überdeckung, konkreter Kamerabewegung, Kameramodell und
STS-Optimierung. Die Maskensegmentierung selbst blieb davon unberührt stabil –
das Rohr wurde in jedem Frame erkannt. Ein isolierter Kausalnachweis für einen
einzelnen Faktor liegt nicht vor.

\subsection{Produktionskonfiguration aus der Matrixauswertung}
\label{sec:produktionskonfiguration}

Auf Grundlage der Matrixauswertung wurde die Produktionskonfiguration auf
\texttt{OPENCV}, 5 FPS und Route A festgelegt. Im Repeat-Lauf erzielen die
\texttt{OPENCV}-A-Läufe bei 720p die höchsten objektmaskierten STS-Metriken
der untersuchten Kameramodelle (29,62\,dB PSNR, 0,902 SSIM und 0,117 LPIPS;
\texttt{SIMPLE\_RADIAL}: 29,43\,dB, 0,897 und 0,122); die
Konfiguration ist über die archivierten Läufe reproduzierbar belegt. Die
Umstellung bleibt dennoch eine datengestützte Projektentscheidung ohne
unabhängige geometrische Referenz: Bei qHD liegt \texttt{SIMPLE\_RADIAL} mit
28,32\,dB vor \texttt{OPENCV} (27,63\,dB), was die nichtmonotone Interaktion
von Auflösung und Kameramodell unterstreicht. Der Standard gilt daher
ausdrücklich für die Produktionsauflösung 720p bei 5 FPS.

Eine nachträgliche Analyse des archivierten 720p-OPENCV-Sparse-Modells ergab
240 registrierte Bilder, 133.723 3D-Punkte, 890.377 Beobachtungen und einen
mittleren Reprojektionsfehler von 0,706\,Pixel. Die entsprechenden Werte lagen
bei \texttt{SIMPLE\_RADIAL} bei 0,707\,Pixel und bei \texttt{PINHOLE} bei
0,705\,Pixel. Alle drei Modelle beschreiben die beobachteten
Bildkorrespondenzen damit fast identisch und mit einem Fehler, der klein
gegenüber der Bildbreite ist; ein absoluter Schwellenwert, ab dem man von
einer „guten" Anpassung sprechen dürfte, existiert jedoch nicht. Der geringe
Unterschied entscheidet daher nicht, welches Modell das reale Objektiv
physikalisch korrekt beschreibt.
Die Reprojektionswerte sind eine Konsistenzprüfung des SfM-Modells und keine
unabhängige Kalibrierung.

Für eine belastbare Aussage über die tatsächliche Verzeichnung wären deshalb
zusätzliche Referenzdaten erforderlich, etwa eine unabhängige
Checkerboard- oder ChArUco-Kalibrierung mit denselben Aufnahmeparametern,
eine Prüfung gerader Kanten vor und nach der Entzerrung oder eine Auswertung
auf zurückgehaltenen Bildkorrespondenzen. Ohne diese Referenz bleibt
\texttt{OPENCV} eine begründete, intern konsistente Produktionsentscheidung
und kein messtechnisch validiertes Kameramodell.

Für die drei Kameramodelle des Repeat-Laufs ergeben sich gemittelte Laufzeiten
von rund 2011 Sekunden (720p), 1650 Sekunden (qHD) und 1199 Sekunden (low).
Diese Werte fließen in die Lauf-Preset-Abfrage
des CLI ein (Kapitel~\ref{sec:implementierung}), mit der zwischen
Bildauflösung und Rechenaufwand gewählt werden kann. Als reine Schrittzeiten
von STS bis Postprocessing ersetzen sie keine vollständige End-to-End-
Zeitmessung.

\subsection{Gemessene End-to-End-Laufzeiten}
\label{sec:laufzeiten-e2e}

Für vollständige End-to-End-Zeiten wurde der Verifikationsbatch
\path{matrix_e2e_verifikation_260826} ausgeführt: drei Matrixläufe der
Produktionskonfiguration (Route A, \texttt{OPENCV}, 5 FPS) über die drei
Auflösungen, jeweils inklusive SAM~3.1, COLMAP, Maskenwarp, Rendering und
Archivierung. Tabelle~\ref{tab:e2e-laufzeiten} fasst die aus dem Batch-Log
rekonstruierten Phasenzeiten zusammen; die Nachlaufphase ist als Obergrenze bis
zum Folgesegmentbeginn ausgewiesen.

\begin{table}[H]
\centering
\small
\begin{tabularx}{\textwidth}{lrrrr}
\toprule
Auflösung & Kopf (SAM3+COLMAP+Warp) & STS--Post & Nachlauf & Gesamt\\
\midrule
720p & 7:49 & 32:23 & $\le$0:52 & $\approx$40:25\\
qHD  & 7:17 & 26:24 & $\le$0:44 & $\approx$33:52\\
low  & 5:27 & 18:11 & --        & $\approx$23:46\\
\bottomrule
\end{tabularx}
\caption{Gemessene End-to-End-Laufzeiten der Produktionskonfiguration
(Verifikationsbatch \path{matrix_e2e_verifikation_260826}, Quelle
\path{e2e_times.csv}). Die Nachlaufphase umfasst Render-, Eval- und
Archivierungsschritte bis zum Folgesegmentbeginn.}
\label{tab:e2e-laufzeiten}
\end{table}

Die Werte bestätigen die Preset-Angabe des CLI und werden durch den
unabhängigen Batch \path{matrix_qualitaetsvergleich_20260818} gestützt, dessen
Segmentlogs nach derselben Methode 40:16, 34:35 und 24:02 Minuten ergeben. Die
Kopfphase skaliert erwartbar mit der Auflösung; die STS--Post-Zeiten dominieren
den Gesamtlauf mit 80--82\,\%. Die timeshare-Angaben begründen keine
Genauigkeitsaussagen und dienen allein der Ressourcenplanung.

\subsection{Vierfeld-Ablation und Produktionsentscheidung}

Die kontrollierte Ablation trennt Original-STS-Gaussians beziehungsweise
SuGaR-Coarse und zwei Tiefenrouten. Tabelle~\ref{tab:mesh-ablation} enthält die
resultierenden Meshgrößen.

\begin{table}[H]
\centering
\small
\begin{tabularx}{\textwidth}{lXrr}
\toprule
Variante & Eingang und Tiefenroute & Vertices & Faces\\
\midrule
A & Original-STS-GS, Diamond-Mesh-Z-Buffer & 68.471 & 135.164\\
B & Original-STS-GS, Gaussian-Rasterizer-Tiefe & 63.266 & 124.976\\
C & SuGaR-Coarse $c=9000$, Diamond-Mesh-Z-Buffer & 79.711 & 156.935\\
D & SuGaR-Coarse $c=9000$, Gaussian-Rasterizer-Tiefe & 68.618 & 135.522\\
\bottomrule
\end{tabularx}
\caption{Ergebnisse der Coarse-only-Vierfeld-Ablation bei identischen
Extraktionsparametern und Seed 42.}
\label{tab:mesh-ablation}
\end{table}

A und B liegen großräumig auf derselben Objektgeometrie, sind aber nicht
identisch. Alle vier Varianten entstammen derselben STS-Basis und teilen
deren lokales Koordinatensystem; der Modellmaßstab ist allerdings nicht gegen
eine reale Referenz geeicht, weshalb die Zentimeterangaben lokale
Modelleinheiten und keine Vermessungswerte sind. Für jeden Variantenvergleich
wurden jeweils 20.000 Punkte gleichmäßig auf einer Oberfläche gesampelt und
gegen die andere Oberfläche gerichtet ausgewertet; daraus ergeben sich pro
Paar zwei gerichtete Mittelwerte (hin und zurück). So lagen A gegen B bei
etwa 1,08 beziehungsweise 0,90\,cm, C gegen A bei etwa 1,79 und 3,39\,cm und
D gegen A bei etwa 1,96 und 2,73\,cm. Seed 42 legt dabei die Zufallsfolgen
von Surface-Sampling und stochastischen Optimierungsschritten fest und macht
die vier Varianten untereinander reproduzierbar vergleichbar. Die größeren
Gegenrichtungsabstände und die visuell stärkere Ausreißerstruktur der
Coarse-Varianten zeigen, dass die zusätzliche Gaussian-Optimierung die
Oberflächenstützung relevant verändert. Diese Werte messen Variantenabstände,
nicht die Genauigkeit zum realen Rohr. Die Ablationsmeshes liegen im
komprimierten Archivband des Ursprungsbatches; Extraktionsparameter und Seed
sind in Tabelle~\ref{tab:versuchsparameter} vollständig dokumentiert, sodass
die Zellen bei Vorliegen des Checkpoints deterministisch reproduzierbar sind.

Route A wird deshalb als Produktionsstandard gewählt: Sie bewahrt die
Original-STS-Gaussians als primäre Geometriehypothese und vermeidet eine
zusätzliche Coarse-Optimierung, die für die Centerline-Anwendung mehr Details,
aber auch mehr Spitzen und Randartefakte erzeugen kann. Route~B verwendet zwar
denselben Original-STS-Eingang, wurde aber nur als Diagnosevariante für die
Tiefenquelle geführt: Die Rasterizer-Tiefe stammt aus der Diskretisierung der
projizierten Gaussians und wirkt an Silhouetten und verdeckten Bereichen
unruhiger, während die projizierte Diamond-Mesh-Tiefe näher an der gelernten
Oberfläche bleibt; zudem liefert B das kleinste Vertex-/Face-Ergebnis der
vier Varianten. Die Wahl bleibt eine begründete Projektentscheidung ohne
unabhängige Ground Truth. Dass die
Coarse-Route dennoch existiert, erklärt sich aus der SuGaR-Startarchitektur:
Sie ist der dort übliche Weg zur Meshgewinnung, während die direkte Extraktion
ohne Coarse-Optimierung erst im Projektverlauf ergänzend erprobt wurde
(vgl. Kapitel~\ref{sec:grundlagen} sowie Abschnitt~\ref{sec:sugar-folgepruefung}).

\subsection{Abgeschlossene SuGaR-Coarse-Folgeprüfung}
\label{sec:sugar-folgepruefung}

Die zwölf Läufe der separaten Matrix
\texttt{matrix\_sugar\_followup\_12} wurden vollständig ausgeführt: je sechs
auf \texttt{SIMPLE\_RADIAL} und \texttt{OPENCV}, bei 2 und 5 FPS sowie in
720p, qHD und Low. Alle zwölf Manifeste tragen den Status \texttt{success};
Masken- und Domänenprüfung, STS-Training, SuGaR-Coarse-Mesh,
SuGaR-Coarse-Rendering, objektmaskierte Metriken und Postprocessing sind damit
für diese zwölf Konfigurationen archiviert. Warum die Coarse-Route neben der
Produktionsroute A weitergeführt wird, ist in Kapitel~\ref{sec:grundlagen}
begründet; Abbildung~\ref{fig:pa-panel-sugar} ordnet die historischen
Entwicklungsstände dieser Route ein.

\begin{figure}[H]
  \centering
  \includegraphics[width=\textwidth]{pa_panel_sugar_iteration.png}
  \caption[Historische SuGaR-Coarse-Zwischenstände]{Historische
  SuGaR-Coarse-Zwischenstände der Vergleichsroute: Meshstufen mit Zielzähler
  9001 sowie bilaterale Vergleichsläufe im 9000er-/10.000er- und
  10.000er-/15.000er-Bereich.}
  \label{fig:pa-panel-sugar}
\end{figure}

Die Mittelwerte der beiden FPS-Profile liegen für
\texttt{SIMPLE\_RADIAL/SuGaR-Coarse} bei 21,38 dB PSNR, 0,764 SSIM und 0,153
LPIPS in 720p; bei qHD bei 21,69 dB, 0,761 und 0,141; bei Low bei 21,89 dB,
0,752 und 0,126. Für \texttt{OPENCV/SuGaR-Coarse} ergeben sich 21,45 dB,
0,764 und 0,153 in 720p; 21,62 dB, 0,764 und 0,137 in qHD sowie 21,81 dB,
0,749 und 0,126 in Low. Tabelle~\ref{tab:sugar-followup-average} fasst die
Mittelwerte zusammen; für ihre Interpretation gilt der in
Kapitel~\ref{sec:grundlagen} eingeführte Ansichtsmetrik-Vorbehalt unverändert.

\begin{table}[H]
\centering
\small
\begin{tabularx}{\textwidth}{l l rrr}
\toprule
Kameramodell & Auflösung & PSNR$\uparrow$ & SSIM$\uparrow$ & LPIPS$\downarrow$\\
\midrule
SIMPLE\_RADIAL & 720p & 21,43 & 0,766 & 0,153\\
SIMPLE\_RADIAL & qHD  & 21,59 & 0,757 & 0,140\\
SIMPLE\_RADIAL & Low  & 21,87 & 0,752 & 0,125\\
OPENCV         & 720p & 21,38 & 0,761 & 0,154\\
OPENCV         & qHD  & 21,60 & 0,754 & 0,144\\
OPENCV         & Low  & 21,84 & 0,747 & 0,126\\
\bottomrule
\end{tabularx}
\caption{SuGaR-Coarse-Mittelwerte aus 2 und 5 FPS der Folgematrix. Die
Werte sind objektmaskierte Ansichtsmetriken und kein Geometrienachweis.}
\label{tab:sugar-followup-average}
\end{table}

Die \texttt{sugar\_coarse\_masked.json}-Dateien enthalten je nach FPS zwölf
beziehungsweise dreißig Evaluationsansichten. Die zugehörigen
\texttt{sugar\_refined\_masked.json}-Dateien sind in allen zwölf Archiven
wegen eines nicht exportierten \texttt{refined.ply} übersprungen. Der neue
Befund weist daher die technische SuGaR-Coarse-Route nach, schließt aber noch
keinen Vergleich der SuGaR-Refined-Stufe ein.

Als Zusatzgrafiken bleiben die gepaarte Differenz zwischen SuGaR-Coarse und
STS sowie die Beziehung zwischen Bildmetrik und archivierter Laufzeit. Sie
sind Diagnose- und Interpretationshilfen, keine neuen Geometriemetriken.

\begin{figure}[H]
  \centering
  \includegraphics[width=\textwidth]{sugar_coarse_vs_sts_delta.pdf}
  \caption{Gepaarte Differenz zwischen SuGaR-Coarse und STS bei identischer
  Follow-up-Konfiguration. Negative PSNR-/SSIM-Werte und positive LPIPS-Werte
  bedeuten in diesen Läufen eine schlechtere SuGaR-Coarse-Bildmetrik; die
  Differenz ist kein Mesh-Genauigkeitsnachweis.}
  \label{fig:pa-sugar-delta}
\end{figure}

\begin{figure}[H]
  \centering
  \includegraphics[width=\textwidth]{metric_vs_runtime.pdf}
  \caption{Explorative Beziehung zwischen objektmaskierten Bildmetriken und
  der archivierten Laufzeit der Pipeline-Segmente STS bis Postprocessing.
  SAM~3.1- und COLMAP-Wandzeit sind in diesen Segmentzeiten nicht enthalten;
  gemessene Gesamtlaufzeiten je Auflösung weist Abschnitt~\ref{sec:laufzeiten-e2e}
  separat aus.}
  \label{fig:pa-metric-runtime}
\end{figure}

In der Laufzeitgrafik ist die x-Achse auf 15--40 Minuten begrenzt und zeigt
die Gesamtlaufzeit des archivierten Pipeline-Laufs. Grün und Lila kodieren STS
und SuGaR-Coarse; kräftige und blasse Farbtöne kodieren 5 beziehungsweise
2 FPS. Kreis, Quadrat und Dreieck stehen für 720p, qHD und Low. Die Kamera
wird über die Achsenbeschriftung angegeben, nicht über eine zusätzliche
SR-/OPENCV-Legende.

\subsection{Funktionsnachweis der Endstufen}

Die erfolgreichen A-Läufe erzeugten neben den STS-Metriken auch Coarse-Ply,
kompatibles OBJ, lokale Centerline, B-Spline und GeoJSON. Die Georeferenzierung
war in den Matrixarchiven mangels validierter CloudCompare-Matrix teilweise ein
Translation-Fallback. Damit ist die technische Kette bis zu lokalen
Geometrieprodukten nachgewiesen; eine globale Genauigkeit ist daraus nicht
ableitbar.

Abbildung~\ref{fig:pa-panel-alurohr} zeigt das für den Testdatensatz erzeugte
Rohrmesh und die daraus abgeleitete lokale Centerline als Ergebnis der
Endstufen.

\begin{figure}[H]
  \centering
  \includegraphics[width=\textwidth]{pa_panel_alurohr_endprodukt.png}
  \caption{Alurohr-Endprodukt: rekonstruiertes Rohrmesh und die daraus
  extrahierte lokale Centerline. Beide Ansichten stammen aus demselben
  Entwicklungsstand (18. Juli 2026).}
  \label{fig:pa-panel-alurohr}
\end{figure}

Beim kanonischen Produktionslauf (\path{Alurohr_THWS.mp4}, OPENCV, 5 FPS,
720p, Route~A) extrahierte DGtal 87 Rohpunkte; der
Grad-10-B-Spline erzeugte 380 Ausgabepunkte. Diese Zahlen belegen die
Ausführbarkeit des Postprocessings. Sie bewerten weder die Lage der Linie im
Rohr noch ihre Formabweichung gegen eine Referenz. Die anschaulichen
Centerline-Bilder der drei Auflösungsstufen finden sich in
Anhang~\ref{app:durchlauf}; dort sind sie als Teil der vollständigen
Durchlaufkette im Kontext der übrigen Zwischenstufen zu sehen.

\subsection{Negative Ergebnisse als Robustheitsbefund}

Mehrere zunächst fehlgeschlagene Läufe lieferten übertragbare Ergebnisse für
die Architektur. Die feste Eval-Liste muss Bildnamen unabhängig von
Dateiendungen abgleichen; die konkrete Eval-Maske muss nichtleer sein; ein
Unterprozess darf nicht die Konfigurations-TSV des Matrixrunners konsumieren;
und SuGaR-Renderings benötigen sowohl den richtigen Importpfad als auch das
erwartete Tensorlayout. Diese Korrekturen sind heute als Quality-Gates oder
explizite Fehlerprüfungen implementiert. Tabelle~\ref{tab:fehlerchronik}
dokumentiert weitere Fälle.
